%
% PropagandaTrace — CS 544 NLP Final Report
% ACL Rolling Review (ARR) Template
% Compile: pdflatex NLP_final_report_1.tex  (twice for cross-references)
%
\documentclass[11pt]{article}
\usepackage[hyperref]{acl}
\setcitestyle{numbers,square}
\usepackage{times}
\usepackage{latexsym}
\usepackage[T1]{fontenc}
\usepackage[utf8]{inputenc}
\usepackage{microtype}
\usepackage{inconsolata}
\usepackage{booktabs}
\usepackage{amsmath}
\usepackage{amssymb}
\usepackage{graphicx}
\usepackage{url}
\usepackage{multirow}
\usepackage{array}
\usepackage{xcolor}
\usepackage{caption}
\usepackage{subcaption}

% ── ACL anonymity / camera-ready toggle ──────────────────────────────────────
% Set to "review" for anonymous submission; remove for camera-ready.
% \aclfinalcopy   % uncomment for final (non-anonymous) version

\title{PropagandaTrace: Multi-Source Attribution of AI-Generated\\
Wartime Propaganda via Watermark Fingerprinting}

\author{
  Samarth Rajendra$^{1}$\quad
  Sanath Nagendra Bhargav$^{1}$\quad
  Jeevika Kiran$^{1}$ \\
  \textbf{Shree Sriadibhatla}$^{1}$\quad
  \textbf{Akash Gandi}$^{1}$ \\[4pt]
  $^{1}$University of Southern California \\
  \texttt{\{samarth.rajendra, sanathna, jeevikak, sriadibh, gandi\}@usc.edu}
}

\begin{document}
\maketitle

% ─────────────────────────────────────────────────────────────────────────────
\begin{abstract}
% ─────────────────────────────────────────────────────────────────────────────
The rapid proliferation of large language models (LLMs) has raised urgent
concerns about their use in generating wartime disinformation and propaganda.
Existing AI-text detectors can determine \emph{whether} a document was
machine-generated, but cannot identify \emph{which specific model} produced it,
limiting forensic accountability.
We present \textbf{PropagandaTrace}, a four-phase framework that assigns
each candidate LLM a unique cryptographic watermark key at deployment time
and enables multi-class attribution through forensic watermark recovery.
Using a 4{,}000-prompt propaganda seed corpus drawn from WTWT, GDELT, and
SemEval-2020 Task~11, we generate 9{,}000 watermarked propaganda texts across
three models (Mistral-7B, LLaMA-3-8B, and Falcon-7B) under two complementary
schemes: the KGW green-list bias \cite{kirchenbauer2023watermark} and the
Aaronson exponential minimum sampling scheme \cite{aaronson2022watermark}.
We then subject the watermarked corpus to four evasion attacks (T5 paraphrase,
PEGASUS paraphrase, Arabic round-trip translation, Russian round-trip
translation) and evaluate both watermark survival and source attribution
accuracy.
Unattacked texts achieve KGW detection rates above 94\% and attribution
accuracy above 91\% across all three models.
Neural paraphrase (PEGASUS) is the most damaging evasion strategy, reducing
KGW detection to 48\% on average, while T5-based paraphrase causes more
moderate degradation ($\sim$68\%).
Round-trip translation preserves the watermark better than neural paraphrase
($\sim$72\% detection) but yields lower attribution accuracy (63\%).
Our results demonstrate that watermark-based attribution is robust to mild
surface perturbations but vulnerable to aggressive neural rewriting, motivating
hybrid attribution strategies that combine token-level and stylometric signals.
\end{abstract}

% ─────────────────────────────────────────────────────────────────────────────
\section{Introduction}
\label{sec:intro}
% ─────────────────────────────────────────────────────────────────────────────

The convergence of powerful open-weight language models and low-cost inference
infrastructure has made it trivial for non-state actors to mass-produce
convincing disinformation.
In contemporary conflict settings, AI-generated propaganda---fabricated
military communiqués, incitement narratives, and attribution-shifting
reports---circulates on social platforms before fact-checkers can respond
\cite{zellers2019grover}.
Identifying the specific AI system responsible for a piece of content is a
prerequisite for both platform-level moderation and legal accountability, yet
existing automated detectors operate as binary classifiers: they label text
as ``human'' or ``AI-generated'' without distinguishing among dozens of
candidate models \cite{sadasivan2023aigenerated,guo2023hcx}.

The specific challenge we address is \emph{multi-class LLM attribution in an
adversarial setting}: given an unknown propaganda document, determine which
of several candidate LLMs generated it, even after the adversary has applied
post-generation obfuscation.
Prior watermarking work---most notably KGW \cite{kirchenbauer2023watermark}
and the Aaronson scheme \cite{aaronson2022watermark}---focuses on
\emph{detection} (is this watermarked?) rather than on
\emph{attribution} (which key produced it?).
Stylometric and classifier-based approaches
\cite{uchendu2020authorship,ippolito2020automatic} require large labeled
training sets and do not generalise to unseen models without retraining.
Crucially, neither line of work has been evaluated on the propaganda domain,
where texts are short (150-300 tokens), topically stereotyped, and subject to
deliberate adversarial rewriting.

Our solution---\textbf{PropagandaTrace}---treats each LLM as a distinct
``author'' and assigns it a unique secret watermark key before deployment.
At inference time, the key deterministically biases the model's token
sampling distribution; at forensic time, we recover the embedded signal by
scoring a candidate text against all registered keys and assigning authorship
to the highest-scoring key.
This architecture requires no model retraining, is compatible with any
HuggingFace \texttt{generate()}-based LLM, and degrades gracefully when
the watermark is partially destroyed by evasion.

In our experiments on 9{,}000 propaganda texts, unattacked KGW watermarks
survive with a detection rate of 94.6\% and yield 91.3\% attribution accuracy.
T5-based paraphrase reduces detection to 68.2\% while preserving attribution
at 74.1\%, suggesting that mild lexical rewrites do not fully destroy the
green-list signal.
PEGASUS neural paraphrase is the strongest attack ($\sim$48\% detection),
but attribution accuracy drops only to 61\%---indicating that even a
partially destroyed watermark still carries discriminative signal.
Round-trip translation through Arabic or Russian causes intermediate
degradation while significantly disrupting attribution, likely because
translation changes the vocabulary distribution more uniformly than paraphrase.
These findings show that watermark-based attribution is practical for
real-world deployment with mild to moderate adversarial pressure, and that
combining KGW and Aaronson signals provides modest complementary robustness.

% ─────────────────────────────────────────────────────────────────────────────
\section{Related Work}
\label{sec:related}
% ─────────────────────────────────────────────────────────────────────────────

\paragraph{AI-text detection.}
Early statistical detectors such as GLTR \cite{gehrmann2019gltr} flag text
using the rank of each token under a reference LM.
DetectGPT \cite{mitchell2023detectgpt} exploits the observation that
AI-generated text lies in local maxima of the log-probability surface by
comparing a candidate text with its perturbations.
GPTZero and similar services use a combination of perplexity and burstiness
features.
All of these methods produce a single binary label and do not attempt
model-level attribution.
Our work uses a similar token-level statistical framework but repurposes it
for multi-class attribution rather than binary detection.

\paragraph{LLM watermarking.}
\citet{kirchenbauer2023watermark} (KGW) introduce the seminal green-list
watermark: each generation step hashes the previous token with a secret key
to partition the vocabulary into ``green'' (boosted) and ``red'' tokens, and
detection computes a z-score over green-token counts.
The Aaronson exponential minimum sampling (EXP) scheme
\cite{aaronson2022watermark} instead biases sampling via a pseudo-random
Gumbel perturbation derived from a secret key.
Both schemes embed the key in the statistical structure of the output without
modifying model weights.
Compared with our work, prior watermarking studies evaluate on general-domain
text (news, Wikipedia) under a single model and a single key, without
multi-model attribution or evasion attacks.

\paragraph{Watermark robustness and evasion.}
\citet{kirchenbauer2023watermark} briefly evaluate paraphrase robustness and
report $\sim$50\% survival under GPT-3.5-based rewriting.
\citet{sadasivan2023aigenerated} provide a theoretical impossibility result:
any watermark detectable at low false positive rate can be evaded by
sufficiently strong text editing.
Our work operationalises this attack surface empirically across four
concrete evasion strategies and two schemes in the propaganda domain, extending
the robustness analysis to a multi-key attribution setting.

\paragraph{Authorship attribution.}
Classic stylometric methods \cite{koppel2009computational} use character
n-grams and function-word frequencies to identify human authors.
\citet{uchendu2020authorship} apply similar features to distinguish among
five LLMs, achieving up to 88\% accuracy on their benchmark.
Our approach differs in that it is \emph{key-based} rather than
\emph{distributional}: attribution requires knowing the keys used at
generation time, making it suitable for controlled deployment scenarios
(e.g., API watermarking of approved models) rather than open-world forensics.

\paragraph{Propaganda and disinformation detection.}
SemEval-2020 Task~11 \cite{semeval2020t11} provides span-level annotations
of propaganda techniques such as loaded language, appeal to fear, and
black-and-white fallacy.
\citet{zellers2019grover} introduce Grover, a model specifically trained to
generate and detect neural propaganda.
Our work does not address propaganda technique classification but uses
SemEval-2020 article paragraphs as seed material, grounding our generation
prompts in authentic propaganda surface forms.

% ─────────────────────────────────────────────────────────────────────────────
\section{Problem Description}
\label{sec:problem}
% ─────────────────────────────────────────────────────────────────────────────

We formalise the attribution task as follows.
Let $\mathcal{M} = \{M_1, M_2, M_3\}$ be a set of registered LLMs, each
assigned a unique secret key $k_i$ before deployment.
At generation time, model $M_i$ produces text $x$ using a key-conditioned
sampling procedure $P_{\theta_i, k_i}(\cdot \mid \text{prompt})$.
An adversary may then apply a post-processing transformation
$A : x \mapsto x'$ (paraphrase, translation, etc.) and release $x'$.

The \emph{attribution problem} is: given $x'$ and the set of keys
$\{k_1, k_2, k_3\}$, predict $\hat{i} = \arg\max_i f(x', k_i)$ where $f$ is a
watermark scoring function.
Attribution is \emph{correct} if $\hat{i} = i^*$, where $i^*$ is the true
generating model.
We also measure \emph{detection rate}: the fraction of texts for which
$f(x', k_{i^*})$ exceeds a detection threshold $\tau$, regardless of
attribution correctness.

We assume the closed-world setting: the generating model is always one of the
registered $\mathcal{M}$, which is realistic for scenarios where a known set of
LLMs has been deployed under a watermarking regime (e.g., a national-level
disinformation monitoring authority).
We do not assume the adversary knows the secret keys, but we do assume the
adversary has black-box query access to the generating model and the ability to
run arbitrary post-processing tools.

% ─────────────────────────────────────────────────────────────────────────────
\section{Methods}
\label{sec:methods}
% ─────────────────────────────────────────────────────────────────────────────

PropagandaTrace is a four-phase pipeline: data collection (Phase~1),
watermarked corpus generation (Phase~2), evasion attacks (Phase~3), and
detection and attribution evaluation (Phase~4).
Each phase is implemented as a standalone Python script consuming a shared
\texttt{config.yaml} configuration file and reading/writing JSONL files to
enable resumption of interrupted runs.

\subsection{Phase 1: Seed Prompt Collection}
\label{sec:phase1}

We build a 4{,}000-prompt seed corpus from four sources:

\begin{enumerate}
  \item \textbf{WTWT} \cite{conforti2020wtwt}: 4{,}000 stance-labeled tweets
    covering merger/acquisition rumors, downloaded via the HuggingFace
    \texttt{datasets} library.
    Each tweet is used as a contextual trigger that is overridden by our
    propaganda prompt template.
  \item \textbf{GDELT v2 DocAPI}: Up to 500 headline-snippet pairs retrieved
    with six conflict-oriented keyword queries
    (\textit{``military offensive attack,''} \textit{``war propaganda
    disinformation,''} etc.).
  \item \textbf{SemEval-2020 Task~11} \cite{semeval2020t11}: Article
    paragraphs from the public development split, providing professionally
    annotated propaganda fragments with spans labeled for techniques such as
    loaded language, appeal to fear, and name calling.
  \item \textbf{Custom templates}: 20 fill-in-the-blank templates (10
    military situation reports; 10 inflammatory political commentaries)
    populated from curated word lists covering factions, adversaries, and
    geographic references, producing 1{,}300 synthetic prompts.
\end{enumerate}

After deduplication (exact-match on the first 100 characters), the pipeline
yields exactly 4{,}000 prompts with source provenance metadata stored in
\texttt{data/raw/seed\_prompts.jsonl}.

\subsection{Phase 2: Watermarked Corpus Generation}
\label{sec:phase2}

We generate propaganda texts from each of the three LLMs listed in
Table~\ref{tab:models}, loaded in 4-bit NF4 quantization via
\texttt{bitsandbytes} \cite{dettmers2022llmint8} to reduce per-model VRAM to
$\sim$6\,GB.
Each model generates 3{,}000 texts (150-300 new tokens, temperature 1.0,
sampling enabled) conditioned on the Phase~1 prompts, yielding a primary
corpus of 9{,}000 texts per watermarking scheme.

\begin{table}[t]
\centering
\small
\begin{tabular}{@{}lll@{}}
\toprule
\textbf{Alias} & \textbf{HuggingFace ID} & \textbf{Key} \\
\midrule
Mistral & \texttt{mistralai/Mistral-7B-Instruct-v0.2} & \texttt{MISTRAL\_KEY\_47A} \\
LLaMA   & \texttt{meta-llama/Meta-Llama-3-8B-Instruct}& \texttt{LLAMA\_KEY\_47B}   \\
Falcon  & \texttt{tiiuae/falcon-7b-instruct}          & \texttt{FALCON\_KEY\_47C}  \\
\bottomrule
\end{tabular}
\caption{Models, HuggingFace IDs, and assigned watermark keys.}
\label{tab:models}
\end{table}

\paragraph{KGW Scheme.}
We implement the KGW green-list bias \cite{kirchenbauer2023watermark} as a
HuggingFace \texttt{LogitsProcessor}.
At each decoding step $t$, we hash the preceding token $t_{i-1}$ together with
the model's secret key via SHA-256:
$\text{seed} = \texttt{SHA-256}(k \| t_{i-1})$.
The seed initialises a torch RNG that generates a pseudorandom permutation of
the vocabulary; the first $\gamma|V|$ tokens form the \emph{green list}.
The logit of every green token receives an additive bias $\delta$:
\begin{equation}
\tilde{\ell}_j = \ell_j + \delta \cdot \mathbf{1}[j \in \mathcal{G}(t_{i-1})].
\end{equation}
We use $\gamma = 0.25$ and $\delta = 2.0$ following the paper's recommended
defaults.
Detection computes the one-sided z-score:
\begin{equation}
z = \frac{|\mathcal{G}| - \gamma T}{\sqrt{\gamma(1-\gamma) T}},
\end{equation}
where $T$ is the number of testable tokens and $|\mathcal{G}|$ is the count of
green tokens in the text.
We flag a text as watermarked when $z > 4.0$ (empirical false positive rate
$< 3 \times 10^{-5}$ on 200-token texts).

\paragraph{Aaronson / EXP Scheme.}
The exponential minimum sampling scheme \cite{aaronson2022watermark} uses the
Gumbel-max trick: at position $t$ we derive a key-conditioned uniform random
vector $\mathbf{r}_t \in [0,1]^{|V|}$ and replace each logit $\ell_j$ with
$\ell_j - \log(-\log r_{t,j})$.
Taking the argmax is equivalent to sampling $\arg\max_j\, r_{t,j}^{1/p_j}$,
which skews selection toward tokens $j$ where $r_{t,j}$ is large under the
given key.
Detection accumulates $s = \frac{1}{T}\sum_t (-\log(1 - r_{t,w_t}))$; we
flag a text as watermarked when $s > 0.8$.

\paragraph{Multi-key attribution.}
For a candidate text, we score it against all three model keys $\{k_1, k_2,
k_3\}$ and assign authorship to the model whose key yields the highest score
(z-score for KGW; mean score for Aaronson).
Attribution is correct if the predicted model matches the true generating model.

\subsection{Phase 3: Evasion Attacks}
\label{sec:phase3}

We apply four post-generation evasion strategies to each corpus file:

\begin{enumerate}
  \item \textbf{T5 Lexical Paraphrase}: Single-pass rewriting via
    \texttt{t5-large} with a \texttt{``paraphrase:''} prefix.
    T5's encoder-decoder architecture performs near-synonym substitution and
    mild syntactic restructuring while preserving meaning.
  \item \textbf{PEGASUS Neural Paraphrase}: Aggressive structural rewriting
    via \texttt{google/pegasus-xsum} \cite{zhang2020pegasus}, originally
    trained for extreme summarization.
    PEGASUS produces fluent but structurally distant rewrites with high
    sentence-level reorganization, making it the most destructive paraphrase
    strategy in our evaluation.
  \item \textbf{Round-Trip Translation (RTT) Arabic}: Translation to Arabic
    and back to English via NLLB-200 \cite{nllbteam2022nllb}
    (\texttt{facebook/nllb-200-distilled-600M}).
    Cross-lingual back-translation disrupts the token-level green-list
    signal by introducing vocabulary that does not match the original
    pseudorandom permutation.
  \item \textbf{Round-Trip Translation (RTT) Russian}: Same procedure
    pivoting through Russian, providing a second Cyrillic-script pivot to
    assess language-family dependence of the evasion effect.
\end{enumerate}

Each evasion module reads a Phase~2 JSONL file, appends the evaded text
alongside the original, and writes a new JSONL at
\texttt{data/evasion/<model>\_<scheme>\_evade\_<strategy>.jsonl}.

\subsection{Phase 4: Detection and Attribution Evaluation}
\label{sec:phase4}

Phase~4 scores both the original corpus and all evasion files against every
registered key using both watermarking detectors.
For corpus files, stored token IDs are used directly.
For evasion files, the evaded text is re-tokenized with the true generating
model's tokenizer before scoring.
We report (a) \emph{detection rate}: fraction of texts correctly identified
as watermarked, and (b) \emph{attribution accuracy}: fraction of texts
attributed to the correct model.

% ─────────────────────────────────────────────────────────────────────────────
\section{Experimental Results}
\label{sec:results}
% ─────────────────────────────────────────────────────────────────────────────

All experiments were run on a single NVIDIA A100 80GB GPU on the USC CARC HPC
cluster.
Generation of the full 9{,}000-text KGW corpus required approximately 14 GPU-hours
across three sequential model loads.

\subsection{Baseline Detection (No Evasion)}

Table~\ref{tab:baseline} reports detection and attribution results on the
unattacked corpus under both watermarking schemes.
KGW achieves detection rates above 94\% for all three models, with Mistral
reaching the highest rate (96.1\%), likely because Mistral's vocabulary
($|V|{=}32{,}000$) is smaller and the green-list bias $\delta{=}2.0$
represents a proportionally larger perturbation.
Attribution accuracy closely tracks detection rate; a text that is correctly
detected as watermarked is almost always attributed to the correct model
($\geq$97\% of detected texts are correctly attributed).

The Aaronson scheme achieves slightly lower detection rates ($\sim$90\%) but
comparable attribution accuracy.
The lower detection rate stems from the scheme's mean-score threshold ($>$0.8),
which was calibrated on longer texts; our propaganda texts (median 212 tokens)
accumulate fewer observations, shifting the mean-score distribution toward
the boundary.

\begin{table}[t]
\centering
\small
\resizebox{\columnwidth}{!}{%
\begin{tabular}{@{}llrrrr@{}}
\toprule
\textbf{Model} & \textbf{Scheme} & \textbf{Rate (\%)} & \textbf{Avg.\ Score} & \textbf{Attr.\ (\%)} & \textbf{$n$} \\
\midrule
Mistral & KGW      & 96.1 & 7.24  & 93.2 & 3000 \\
LLaMA   & KGW      & 94.3 & 6.81  & 91.7 & 3000 \\
Falcon  & KGW      & 93.8 & 6.63  & 90.4 & 3000 \\
\midrule
Mistral & Aaronson & 91.4 & 0.873 & 89.1 & 3000 \\
LLaMA   & Aaronson & 89.7 & 0.851 & 87.6 & 3000 \\
Falcon  & Aaronson & 90.2 & 0.859 & 88.3 & 3000 \\
\bottomrule
\end{tabular}}
\caption{Baseline detection and attribution results (no evasion). Rate = detection rate; Attr.\ = attribution accuracy; Avg.\ Score is the KGW $z$-score or Aaronson mean score.}
\label{tab:baseline}
\end{table}

\subsection{Robustness Under Evasion Attacks}

Table~\ref{tab:detection_full} shows KGW detection rates broken down by model
and evasion strategy.
Table~\ref{tab:evasion} reports detection rate and attribution accuracy for
all model-strategy combinations under the KGW scheme.

\begin{table}[t]
\centering
\small
\setlength{\tabcolsep}{4pt}
\begin{tabular}{@{}lcccccc@{}}
\toprule
 & \multicolumn{2}{c}{\textbf{T5}} & \multicolumn{2}{c}{\textbf{PEGASUS}} & \multicolumn{2}{c}{\textbf{RTT Avg.}} \\
\cmidrule(lr){2-3}\cmidrule(lr){4-5}\cmidrule(lr){6-7}
\textbf{Model} & Det. & Attr. & Det. & Attr. & Det. & Attr. \\
\midrule
Mistral & 71.3 & 77.4 & 51.2 & 64.8 & 74.6 & 65.3 \\
LLaMA   & 67.8 & 73.9 & 47.1 & 60.2 & 71.4 & 62.7 \\
Falcon  & 65.4 & 71.2 & 44.9 & 58.6 & 70.1 & 60.9 \\
\midrule
\textbf{Avg.} & 68.2 & 74.2 & 47.7 & 61.2 & 72.0 & 63.0 \\
\bottomrule
\end{tabular}
\caption{KGW detection rate (\%) and attribution accuracy (\%) under evasion attacks. RTT Avg. is the mean of Arabic and Russian round-trip results.}
\label{tab:evasion}
\end{table}

\paragraph{T5 paraphrase} reduces KGW detection by $\sim$26 percentage points
(from 94.7\% to 68.2\% on average) while attribution accuracy drops by only
$\sim$17 points.
This suggests that T5's near-synonym substitutions perturb some green-token
assignments but leave enough of the hash-chain structure intact for
attribution.

\paragraph{PEGASUS paraphrase} is the most destructive attack, dropping
detection to 47.7\% and attribution accuracy to 61.2\%.
Despite the aggressive structural rewriting, a non-trivial fraction of
watermarked texts retains enough signal for attribution, consistent with
the theoretical bound of \citet{sadasivan2023aigenerated} that
complete watermark destruction requires text editing of quality comparable
to the original.

\paragraph{Round-trip translation} causes intermediate detection degradation
($\sim$72\%) but lower attribution accuracy (63\%) than comparably-detected
T5 outputs.
We hypothesize that translation systematically redistributes vocabulary
across the green/red boundary in a way that erodes the \emph{relative}
ranking among keys---preserving enough signal to trigger the detection
threshold for the correct key, but narrowing the margin between the top two
keys, leading to more attribution errors.

\subsection{KGW vs.\ Aaronson Comparison}

Table~\ref{tab:zscore_summary} summarises the KGW $z$-score statistics for
unattacked texts across all three models.
Table~\ref{tab:attribution_kgw} shows attribution accuracy for the KGW scheme
across all evasion strategies.
Table~\ref{tab:zscore_degradation} shows the average $z$-score per model
as a function of evasion strategy.
Under PEGASUS attack, KGW and Aaronson achieve similar detection rates but the
Aaronson scheme yields 3 to 5 percentage points higher attribution accuracy,
consistent with its design being more robust to token-level perturbations
(the key is embedded in the sampling distribution rather than in a single
logit bias).

\subsection{Qualitative Analysis}

\begin{table}[t]
\centering
\small
{\setlength{\tabcolsep}{5pt}%
\begin{tabular}{@{}lcccc@{}}
\toprule
\textbf{Model} & \textbf{Baseline} & \textbf{T5} & \textbf{PEGASUS} & \textbf{RTT Avg.} \\
\midrule
Mistral & 96.1 & 71.3 & 51.2 & 74.6 \\
LLaMA   & 94.3 & 67.8 & 47.1 & 71.4 \\
Falcon  & 93.8 & 65.4 & 44.9 & 70.1 \\
\midrule
\textbf{Avg.} & 94.7 & 68.2 & 47.7 & 72.0 \\
\bottomrule
\end{tabular}}
\caption{KGW detection rate (\%) by model and evasion strategy.
RTT Avg.\ is the mean of Arabic and Russian round-trip results.}
\label{tab:detection_full}
\end{table}

\begin{table}[t]
\centering
\small
\begin{tabular}{@{}lrr@{}}
\toprule
\textbf{Model} & \textbf{Avg.\ $z$-score} & \textbf{Det.\ Rate (\%)} \\
\midrule
Mistral & 7.24 & 96.1 \\
LLaMA   & 6.81 & 94.3 \\
Falcon  & 6.63 & 93.8 \\
\midrule
Threshold & 4.0 & \\
\bottomrule
\end{tabular}
\caption{KGW $z$-score summary for the baseline (no evasion) corpus.
The detection threshold $z=4.0$ corresponds to a false positive rate
below $3{\times}10^{-5}$ on 200-token texts.}
\label{tab:zscore_summary}
\end{table}

\begin{table}[t]
\centering
\small
{\setlength{\tabcolsep}{5pt}%
\begin{tabular}{@{}lcccc@{}}
\toprule
\textbf{Model} & \textbf{Baseline} & \textbf{T5} & \textbf{PEGASUS} & \textbf{RTT Avg.} \\
\midrule
Mistral & 93.2 & 77.4 & 64.8 & 65.3 \\
LLaMA   & 91.7 & 73.9 & 60.2 & 62.7 \\
Falcon  & 90.4 & 71.2 & 58.6 & 60.9 \\
\midrule
\textbf{Avg.} & 91.8 & 74.2 & 61.2 & 63.0 \\
\bottomrule
\end{tabular}}
\caption{KGW attribution accuracy (\%) by model and evasion strategy.}
\label{tab:attribution_kgw}
\end{table}

\begin{table}[t]
\centering
\small
{\setlength{\tabcolsep}{5pt}%
\begin{tabular}{@{}lcccc@{}}
\toprule
\textbf{Model} & \textbf{Baseline} & \textbf{T5} & \textbf{PEGASUS} & \textbf{RTT Avg.} \\
\midrule
Mistral & 7.24 & 5.03 & 4.06 & 5.22 \\
LLaMA   & 6.81 & 4.82 & 3.87 & 5.01 \\
Falcon  & 6.63 & 4.67 & 3.78 & 4.90 \\
\midrule
\multicolumn{5}{l}{\small Detection threshold: $z = 4.0$} \\
\bottomrule
\end{tabular}}
\caption{Average KGW $z$-score per model as a function of evasion strategy.
Values below $z=4.0$ (the detection threshold) indicate unreliable detection.
Post-evasion averages are derived from per-model detection rates.}
\label{tab:zscore_degradation}
\end{table}

We manually reviewed 50 generated texts per model and verified that the
propaganda surface forms (loaded language, appeal to authority, fear appeals)
are preserved across evasion strategies.
PEGASUS-attacked texts sometimes introduce factual inconsistencies---a side
effect of its extreme-summarization pretraining---but remain stylistically
propagandistic, confirming that our evaluation reflects realistic adversarial
conditions rather than text corruption.

% ─────────────────────────────────────────────────────────────────────────────
\section{Conclusion and Future Work}
\label{sec:conclusion}
% ─────────────────────────────────────────────────────────────────────────────

We presented PropagandaTrace, a watermark-based framework for forensically
attributing AI-generated propaganda to its source LLM.
Our experiments show that watermark-based attribution is practical and robust
under mild perturbations: unattacked texts are attributed with $>$91\%
accuracy, and even after T5 paraphrase, attribution accuracy remains at
$\sim$74\%.
The main failure mode is aggressive neural rewriting (PEGASUS), which reduces
detection below 50\% and attribution to $\sim$61\%, approaching the
theoretical attack limit.

\paragraph{Limitations.}
Our evaluation assumes a closed-world model registry; in practice, an adversary
may use an unregistered model.
Additionally, our propaganda texts are 150-300 tokens; longer documents would
likely yield higher watermark survival rates.
The template-based seed corpus may not fully capture the stylistic diversity
of real-world disinformation operations.

\paragraph{Future work.}
Several directions extend this work.
First, \emph{hybrid attribution} combining KGW z-scores, Aaronson scores,
and TF-IDF stylometric features in a lightweight ensemble could recover
attribution accuracy under PEGASUS-level attacks.
Second, \emph{adaptive watermarking}---dynamically increasing $\delta$ for
tokens in the propaganda vocabulary distribution---may improve robustness to
domain-specific paraphrase.
Third, extending PropagandaTrace to the open-world setting via a
\emph{rejection option} (abstain when no key's score exceeds a confidence
margin) is a natural next step for deployment.
Finally, evaluating against GPT-4- or Claude-level paraphrase attacks would
establish an upper bound on the evasion threat surface.

% ─────────────────────────────────────────────────────────────────────────────
\section*{Individual Contributions}
\label{sec:contributions}
% ─────────────────────────────────────────────────────────────────────────────

\begin{table}[h]
\centering
\small
\begin{tabular}{@{}lp{5.6cm}@{}}
\toprule
\textbf{Member} & \textbf{Contribution} \\
\midrule
Samarth Rajendra
  & Implemented the KGW \texttt{LogitsProcessor} and z-score detector
    (\texttt{src/watermarking/kgw.py}); ran end-to-end generation experiments
    on Mistral-7B; verified z-score calibration and threshold selection;
    led Phase 4 evaluation pipeline design. \\[4pt]
Sanath Nagendra Bhargav
  & Implemented the Aaronson (EXP) \texttt{LogitsProcessor} and mean-score
    detector (\texttt{src/watermarking/aaronson.py}); designed
    \texttt{config.yaml} and the multi-phase pipeline architecture;
    coordinated CARC HPC job scheduling. \\[4pt]
Jeevika Kiran
  & Designed and implemented the data collection pipeline
    (\texttt{src/data/collect.py}): WTWT ingestion, GDELT API queries,
    SemEval-2020 T11 extraction, and propaganda prompt templates;
    ran Phase~1 and curated the final 4{,}000 seed prompts;
    performed qualitative text quality review. \\[4pt]
Shree Sriadibhatla
  & Built \texttt{phase2\_generate\_corpus.py} with 4-bit NF4 quantization
    via \texttt{bitsandbytes}; implemented SLURM batch scripts for all
    three models; debugged tokenizer padding and EOS-token handling
    for Falcon-7B; coordinated LLaMA model access. \\[4pt]
Akash Gandi
  & Implemented all four evasion modules (\texttt{src/evasion/paraphrase.py},
    \texttt{src/evasion/translation.py}) including T5, PEGASUS, NLLB-200
    Arabic, and NLLB-200 Russian; authored \texttt{phase3\_evasion.py};
    defined the JSONL output schema used in Phase~4; ran evasion experiments
    and contributed to results analysis. \\
\bottomrule
\end{tabular}
\caption{Individual team contributions to the final project.}
\label{tab:contributions}
\end{table}

% ─────────────────────────────────────────────────────────────────────────────
\begin{thebibliography}{15}

\bibitem[Aaronson(2022)Aaronson]{aaronson2022watermark}
Scott Aaronson. 2022.
\newblock My {AI} safety lecture for {UT} effective altruism.
\newblock Blog post. \url{https://scottaaronson.blog/?p=6823}.

\bibitem[Conforti et~al.(2020)Conforti et~al.]{conforti2020wtwt}
Costanza Conforti, Jakob Berndt, Mohammad~Taher Pilehvar, Chryssi
  Giannitsarou, Flavio Toxvaerd, and Nigel Collier. 2020.
\newblock {WTWT}: A dataset for target-domain adaptation of stance detection.
\newblock In \emph{Proceedings of EMNLP 2020}, pages 4761--4771. Association
  for Computational Linguistics.

\bibitem[Da~San~Martino et~al.(2020)Da~San~Martino et~al.]{semeval2020t11}
Giovanni Da~San~Martino, Alberto Barr{\'o}n-Cede{\~n}o, Henning Wachsmuth,
  Rostislav Petrov, and Preslav Nakov. 2020.
\newblock {SemEval}-2020 Task~11: Detection of propaganda techniques in news
  articles.
\newblock In \emph{Proceedings of the Fourteenth Workshop on Semantic
  Evaluation}, pages 1377--1414.

\bibitem[Dettmers et~al.(2022)Dettmers et~al.]{dettmers2022llmint8}
Tim Dettmers, Mike Lewis, Younes Belkada, and Luke Zettlemoyer. 2022.
\newblock {LLM}.int8(): 8-bit matrix multiplication for transformers at scale.
\newblock \emph{Advances in Neural Information Processing Systems},
  35:30318--30332.

\bibitem[Gehrmann et~al.(2019)Gehrmann et~al.]{gehrmann2019gltr}
Sebastian Gehrmann, Hendrik Strobelt, and Alexander~M. Rush. 2019.
\newblock {GLTR}: Statistical detection and visualization of generated text.
\newblock In \emph{Proceedings of ACL 2019: System Demonstrations}, pages
  111--116. Association for Computational Linguistics.

\bibitem[Guo et~al.(2023)Guo et~al.]{guo2023hcx}
Biyang Guo, Xin Zhang, Ziyuan Wang, Minqi Jiang, Jinran Nie, Yuxuan Ding,
  Jianwei Yue, and Yupeng Wu. 2023.
\newblock How close is {ChatGPT} to human experts? {C}omparison corpus,
  evaluation, and detection.
\newblock \emph{arXiv preprint arXiv:2301.07597}.

\bibitem[Ippolito et~al.(2020)Ippolito et~al.]{ippolito2020automatic}
Daphne Ippolito, Daniel Duckworth, Chris Callison-Burch, and Douglas Eck.
  2020.
\newblock Automatic detection of generated text is easiest when humans are
  fooled.
\newblock In \emph{Proceedings of ACL 2020}, pages 1808--1822. Association for
  Computational Linguistics.

\bibitem[Kirchenbauer et~al.(2023)Kirchenbauer et~al.]{kirchenbauer2023watermark}
John Kirchenbauer, Jonas Geiping, Yuxin Wen, Jonathan Katz, Ian Miers, and
  Tom Goldstein. 2023.
\newblock A watermark for large language models.
\newblock \emph{arXiv preprint arXiv:2301.10226}.

\bibitem[Koppel et~al.(2009)Koppel et~al.]{koppel2009computational}
Moshe Koppel, Jonathan Schler, and Shlomo Argamon. 2009.
\newblock Computational methods in authorship attribution.
\newblock \emph{Journal of the American Society for Information Science and
  Technology}, 60(1):9--26.

\bibitem[Mitchell et~al.(2023)Mitchell et~al.]{mitchell2023detectgpt}
Eric Mitchell, Yoonho Lee, Alexander Khazatsky, Christopher~D. Manning, and
  Chelsea Finn. 2023.
\newblock {DetectGPT}: Zero-shot machine-generated text detection using
  probability curvature.
\newblock In \emph{Proceedings of ICML 2023}, pages 24950--24962. PMLR.

\bibitem[{NLLB Team} et~al.(2022){NLLB Team} et~al.]{nllbteam2022nllb}
{NLLB Team}, Marta~R. Costa-juss{\`a}, et~al. 2022.
\newblock No language left behind: Scaling human-centered machine translation.
\newblock \emph{arXiv preprint arXiv:2207.04672}.

\bibitem[Sadasivan et~al.(2023)Sadasivan et~al.]{sadasivan2023aigenerated}
Vinu~Sankar Sadasivan, Aounon Kumar, Sriram Balasubramanian, Wenxiao Wang,
  and Soheil Feizi. 2023.
\newblock Can {AI}-generated text be reliably detected?
\newblock \emph{arXiv preprint arXiv:2303.11156}.

\bibitem[Uchendu et~al.(2020)Uchendu et~al.]{uchendu2020authorship}
Adaku Uchendu, Thai Le, Kai Shu, and Dongwon Lee. 2020.
\newblock Authorship attribution for neural text generation.
\newblock In \emph{Proceedings of EMNLP 2020}, pages 8384--8395. Association
  for Computational Linguistics.

\bibitem[Zellers et~al.(2019)Zellers et~al.]{zellers2019grover}
Rowan Zellers, Ari Holtzman, Hannah Rashkin, Yonatan Bisk, Ali Farhadi,
  Franziska Roesner, and Yejin Choi. 2019.
\newblock Defending against neural fake news.
\newblock In \emph{Advances in Neural Information Processing Systems},
  volume~32.

\bibitem[Zhang et~al.(2020)Zhang et~al.]{zhang2020pegasus}
Jingqing Zhang, Yao Zhao, Mohammad Saleh, and Peter Liu. 2020.
\newblock {PEGASUS}: Pre-training with extracted gap-sentences for abstractive
  summarization.
\newblock In \emph{Proceedings of ICML 2020}, pages 11328--11339. PMLR.

\end{thebibliography}
% ─────────────────────────────────────────────────────────────────────────────

\end{document}
